% !TeX root = ../main.tex
\chapter{Metric Equations}\label{appendix:metric_equations}

This appendix provides the formal definitions of lexical metrics referenced in Chapter~\ref{chapter:experiments}.

\section{BLEU}
For candidate report $c$ and references $\{r\}$:
\begin{equation}
\mathrm{BLEU} = \mathrm{BP}\cdot \exp\!\left(\sum_{n=1}^{N} w_n \log p_n\right),
\end{equation}
where $p_n$ is modified n-gram precision and $w_n$ are n-gram weights. The brevity penalty is:
\begin{equation}
\mathrm{BP}=\begin{cases}
1 & c_{len} > r_{len}^* \\
\exp\left(1-\frac{r_{len}^*}{c_{len}}\right) & \text{otherwise.}
\end{cases}
\end{equation}

\section{ROUGE-L}
Using longest common subsequence (LCS):
\[
R_{LCS}=\frac{\mathrm{LCS}(c,r)}{|r|},\qquad
P_{LCS}=\frac{\mathrm{LCS}(c,r)}{|c|},
\]
\begin{equation}
F_{LCS}=\frac{(1+\beta^2)R_{LCS}P_{LCS}}{R_{LCS}+\beta^2P_{LCS}}.
\end{equation}

\section{METEOR}
\begin{equation}
F_{\mathrm{mean}}=\frac{(1+\alpha)PR}{R+\alpha P},
\end{equation}
\begin{equation}
\mathrm{METEOR}=F_{\mathrm{mean}}\cdot (1-\mathrm{Penalty}).
\end{equation}

\section{CIDEr}
With TF-IDF-weighted n-gram vectors:
\begin{equation}
\mathrm{CIDEr}_n(c,r)=\frac{1}{M}\sum_{i=1}^{M}
\frac{\mathbf{g}^n(c)\cdot\mathbf{g}^n(r_i)}{\|\mathbf{g}^n(c)\|\|\mathbf{g}^n(r_i)\|}.
\end{equation}
